\documentclass[11pt]{article}

\usepackage[margin=1in]{geometry}
\usepackage{fancyhdr}

\title{\Large PLUS Session Outline \\ \normalsize \textit{MATH 1554:} Linear Algebra}
\author{Ahmed Mohammed}
\date{29th of August, 2023}

% ------- Set up header -------
\renewcommand{\sectionmark}[1]{\markright{#1}} % from https://tex.stackexchange.com/a/317393/227665
\pagestyle{fancy}
\fancyhead[L]{29th of August, 2023}
% \fancyhead[C]{\textit{\nouppercase{\rightmark}}}
\fancyhead[C]{\textit{MATH 1554 PLUS Session Outline}}
\fancyhead[R]{Mohammed \textbf{\thepage}}
% -------------------------------------------------------

\begin{document}
%\maketitle
\thispagestyle{fancy}
\section{\centering Open Activity}
	\begin{enumerate}
		\item \textit{Activity}. I will provide some questions on the slide that concerns definitions as a way to revise concepts covered in the past week and a half.
						I will provide the students some time to respond to the questions asked, and each group will elect a person to answer in behalf
						of that group.
		\item \textit{Collaboration}. Each member of a specific group will be discussing the definition amongst each other, so they'll have to work together
						to reach to an answer.
		\item \textit{Intended Goal}. The purpose of this activity is to serve as a revision for the concepts covered in previous lectures. Also, it'll help me
						understand where the students stand so that I may focus on some questions in the worksheet.
	\end{enumerate}

\section{\centering Main Activity}
	\begin{enumerate}
		\item \textit{Activity}. Students will be given a worksheet that goes over the previous 4 sections covered in lecture. Each group will work on an "activity",
						or rather a multi-part question. Once the allocated time is done, each group will elect a person to tell the whole session the answer
						they reached.
		\item \textit{Collaboration}. As students answer the questions, they'll have to work together. 
		\item \textit{Intended Goal}. The purpose of this activity is to test the students' understanding of concepts.
	\end{enumerate}

\section{\centering Closing Activity}
	\begin{enumerate}
		\item \textit{Activity}. In this activity, students will be asked two questions that foreshadow the next topic: homogenous and inhomogenuous equations,
			and how homogenous equations help us get closer to the solutions of the corresponding consistent inhomogenous system of linear equations.
		\item \textit{Collaboration}. \textsc{N/A}
		\item \textit{Intended Goal}. This activity will help students prepare for the next lecture and session.
	\end{enumerate}
\end{document}